\section{A nonperturbative free-sea dielectric estimate}
\label{fd-section}

This section gives an independent quadratic-field estimate.  It holds on
the entire finite electron--positron Fock space, not just on quasi-free
states or on the neutral sector.  It is useful for distinguishing a proved
free-sea response bound from an assumption about an interacting vacuum.

Let $\mathbb T_L^3$ have volume $V=L^3$, and let $P_f$ project onto the
momenta $p\in(2\pi/L)\mathbb Z^3$ with $|p|\le\Lambda$.  On the resulting
four-spinor space set
\[
 D_0(p)=\boldsymbol{\mathsf a}\cdot p+m\mathsf b,
 \quad E(p)=\sqrt{|p|^2+m^2},\quad
 P_\pm(p)=\frac12\left(1\pm\frac{D_0(p)}{E(p)}\right),\qquad m>0.
\]
The positive particle--hole Hamiltonian and the normal-ordered scalar
coupling are
\[
 H_D=d\Gamma(D_0)-\operatorname{Tr}(P_-D_0),\qquad
 n(\phi)=d\Gamma(\Phi)-\operatorname{Tr}(P_-\Phi),\qquad
 \Phi=P_f\phi P_f.
\]
Here $\phi\in H^1(\mathbb T_L^3)\cap L^\infty(\mathbb T_L^3)$ is real.
The sign of a physical electric potential can be included in $\phi$.
Use the Fourier convention
$\widehat\phi(k)=\int_{\mathbb T_L^3}e^{-ik\cdot x}\phi(x)\,dx$, so that
$\|\nabla\phi\|_2^2=V^{-1}\sum_k|k|^2|\widehat\phi(k)|^2$.

\subsection{The cutoff off-diagonal block and its gradient bounds}

Put $B=P_+\Phi P_-$ and define the weighted squared norm
\begin{equation}
 \mathcal S_{\Lambda,L}(\phi)
 :=\sum_{i\in+,\,j\in-}
       \frac{|B_{ij}|^2}{E_i+|E_j|},
 \label{fd-weighted-norm}
\end{equation}
where the matrix entries are taken in free Dirac eigenbases.  The following
bounds retain the sharp fermion cutoff:
\begin{align}
 \|B\|_{\mathrm{HS}}^2
 &\le C^{(2)}_{\Lambda,L}\|\nabla\phi\|_2^2,
 &C^{(2)}_{\Lambda,L}
 &=\frac1{2V}\sum_{|p|\le\Lambda}E(p)^{-2},
 \label{fd-hs-bound}\\
 \mathcal S_{\Lambda,L}(\phi)
 &\le C^{(3)}_{\Lambda,L}\|\nabla\phi\|_2^2,
 &C^{(3)}_{\Lambda,L}
 &=\frac1{4V}\sum_{|p|\le\Lambda}E(p)^{-3}.
 \label{fd-weighted-bound}
\end{align}
In particular, for every $L\ge L_0>0$ one can use the volume-independent
constants
\begin{align}
 \overline C^{(2)}_{\Lambda,L_0}
 &=\frac1{2m^2L_0^3}+\frac{13\Lambda}{8\pi^3},
 \label{fd-c2}\\
 \overline C^{(3)}_{\Lambda,L_0}
 &=\frac1{4m^3L_0^3}
       +\frac{13}{16\pi^3}\bigl(2+\log_+(\Lambda/m)\bigr),
 \qquad\log_+x=\max\{\log x,0\}.
 \label{fd-c3}
\end{align}

\paragraph{Proof of the gradient bounds.}
The Dirac trace identity for the scalar coupling is
\begin{align}
 \operatorname{tr}_{\mathbb C^4}
       \bigl(P_+(p+k)P_-(p)\bigr)
 &=1-\frac{p\cdot(p+k)+m^2}{E(p)E(p+k)}\nonumber\\
 &=\frac{|k|^2-(E(p+k)-E(p))^2}{2E(p)E(p+k)}
 \le\frac{|k|^2}{2E(p)E(p+k)}.
 \label{fd-projector-trace}
\end{align}
For $a,b>0$,
\[
 \frac1{ab}\le\frac12(a^{-2}+b^{-2}),\qquad
 \frac1{ab(a+b)}\le\frac14(a^{-3}+b^{-3}).
\]
The second inequality follows from
$(a+b)(a^3+b^3)\ge4a^2b^2$.
Insert \eqref{fd-projector-trace} in the Fourier expressions for
$\|B\|_{\mathrm{HS}}^2$ and \eqref{fd-weighted-norm}, and use these two
scalar inequalities.  The sums over $p$ and $p+k$, both required to belong
to the cutoff band, can each be bounded by the full band sum.  Parseval
then proves \eqref{fd-hs-bound}--\eqref{fd-weighted-bound}.  Equivalently,
\[
 \mathcal S_{\Lambda,L}(\phi)
 =\frac1V\sum_k\Pi_{\Lambda,L}(k)|\widehat\phi(k)|^2,
 \quad
 \Pi_{\Lambda,L}(k)=\frac1V
 \sum_{\substack{p,p+k\ \mathrm{in\ the\ band}}}
 \frac{\operatorname{tr}(P_+(p+k)P_-(p))}{E(p+k)+E(p)},
\]
and $0\le\Pi_{\Lambda,L}(k)\le C^{(3)}_{\Lambda,L}|k|^2$.
In particular its zero mode vanishes exactly.

For completeness, the lattice bounds in
\eqref{fd-c2}--\eqref{fd-c3} can be checked without a continuum
approximation.  Set $h=2\pi/L$, $a=m/h$, and $J=\lfloor\Lambda/h\rfloor$.
The shell $|n|_\infty=j$ has $24j^2+2\le26j^2$ integer points.
For the inverse-square sum, this gives
\[
 \frac1V\sum_{0<|p|\le\Lambda}E(p)^{-2}
 \le\frac{26}{Vh^2}\sum_{j=1}^J\frac{j^2}{a^2+j^2}
 \le\frac{26\Lambda}{8\pi^3}.
\]
For the inverse-cube sum, split at $j=a$.  The part $j\le a$ is bounded
by $1$ using $\sum_{j=1}^n j^2=n(n+1)(2n+1)/6$, and the remaining harmonic
sum is bounded by $1+\log_+(\Lambda/m)$.  Consequently
\[
 \frac1V\sum_{0<|p|\le\Lambda}E(p)^{-3}
 \le\frac{26}{8\pi^3}\bigl(2+\log_+(\Lambda/m)\bigr).
\]
Adding the $p=0$ term proves the stated constants.  If $a<1$, the first
part is empty and the same bound follows from
$\sum_{j=1}^Jj^{-1}\le1+\log J$; if $J=0$, the sum is empty.

These are bounds for the compressed multiplication operator $\Phi$ inside
the cutoff one-particle space.  In particular
$\|[P_-,\Phi]\|_{\mathrm{HS}}^2=2\|B\|_{\mathrm{HS}}^2$ there.
They do not assert the same formula for a commutator with a projection
acting on unretained modes.  No cutoff-boundary term has been discarded:
restricting both retained momentum indices only decreases the positive
sums used in the proof.

\subsection{A nonperturbative operator inequality on the full Fock space}

\paragraph{Proposition.}
Let $\tau>0$ and $v=\|\phi\|_\infty<\tau m$.  Then
\begin{align}
 \tau H_D+n(\phi)
 &\ge-\sum_{i\in+,\,j\in-}
       \frac{|B_{ij}|^2}{\tau(E_i+|E_j|)-2v}\,I\nonumber\\
 &\ge-\frac{\mathcal S_{\Lambda,L}(\phi)}{\tau-v/m}\,I
 \ge-\frac{\overline C^{(3)}_{\Lambda,L_0}}
                 {\tau-v/m}\|\nabla\phi\|_2^2\,I
 \quad(L\ge L_0).
 \label{fd-dielectric}
\end{align}
The constants are independent of the charge sector and of the number of
real or virtual pairs.  This is an exact inequality, not a second-order
expansion in the potential.

\paragraph{Proof.}
For any normalized many-body state its one-particle density matrix obeys
$0\le\gamma\le1$ by the CAR, whether or not the state is quasi-free.
Put $Q=\gamma-P_-$ and
$K=Q_{++}-Q_{--}\ge0$.  The inequality $\gamma^2\le\gamma$ implies
\begin{equation}
 Q^2\le K,
 \qquad T:=\operatorname{Tr}(D_0Q)
       =\operatorname{Tr}(|D_0|K).
 \label{fd-q-square}
\end{equation}
The diagonal potential terms are bounded below by
$-v\operatorname{Tr}K$.  Since $\tau|D_0|-v$ is positive,
\[
 \tau T-v\operatorname{Tr}K
 \ge\operatorname{Tr}\bigl((\tau|D_0|-v)Q^2\bigr)
 \ge\sum_{i\in+,\,j\in-}
       [\tau(E_i+|E_j|)-2v]|Q_{ij}|^2.
\]
The omitted diagonal-block squares in the last inequality are positive.
The off-diagonal potential contribution is
$2\operatorname{Re}\sum_{i,j}\overline{B_{ij}}Q_{ij}$.
Completing each scalar square proves the first line of
\eqref{fd-dielectric}.  The second uses
$\tau(E_i+|E_j|)-2v\ge(\tau-v/m)(E_i+|E_j|)$,
and the final bound is \eqref{fd-weighted-bound}.
Since every normalized state satisfies the estimate, it is an operator
inequality on the entire finite Fock space.

\paragraph{Constant shifts and the threshold condition.}
In the exact charge sector $N_e-N_p=q$, a constant shift gives
$n(\phi)=n(\phi-c)+cq$.  Thus one may replace $v$ by
$\|\phi-c\|_\infty$ in \eqref{fd-dielectric} and add $cq$ to its right
side.  On the neutral sector this allows centering the potential by any
constant.  A threshold restriction cannot simply be omitted from an
all-sector gradient-only bound: if $\phi=-c$ with $c>\tau m$, its gradient
vanishes, but a zero-momentum electron has energy $\tau m-c<0$.

There is also a precise estimate with an explicit threshold penalty.
Let $0<\sigma<\tau$, denote the positive operators $|D_0|$ on the two
spectral subspaces by $E_\pm$, and write $A_-:=\max\{-A,0\}$ for the
negative part of a self-adjoint matrix.  Then
\begin{multline}
 \tau H_D+n(\phi)\ge
 -\frac{\mathcal S_{\Lambda,L}(\phi)}{\sigma}
 -\operatorname{Tr}\bigl[((\tau-\sigma)E_++\Phi_{++})_-\bigr]\\
 -\operatorname{Tr}\bigl[((\tau-\sigma)E_- -\Phi_{--})_-\bigr].
 \label{fd-threshold}
\end{multline}
To prove this, use $\sigma T$ and the off-diagonal potential to complete
the weighted squares with denominator $\sigma(E_i+|E_j|)$.
The remaining diagonal expectation is bounded below by the two displayed
negative traces because the electron and hole diagonal density matrices
each lie between zero and one.  No assertion of volume independence is
made for this threshold penalty when the gap condition fails.

\subsection{A precisely scoped mean-field consequence}

Suppose a separately established estimate for a photon-free mean-field
functional gives
\begin{equation}
 \mathcal E_0(\gamma)\ge\theta\operatorname{Tr}(D_0Q)
 \quad\text{for all }0\le\gamma\le1,
 \qquad\theta>0.
 \label{fd-hf-hypothesis}
\end{equation}
Then, for $\|\phi\|_\infty<\theta m$,
\begin{equation}
 \mathcal E_0(\gamma)+\operatorname{Tr}(\Phi Q)
 \ge-\frac{\overline C^{(3)}_{\Lambda,L_0}}
                 {\theta-\|\phi\|_\infty/m}
       \|\nabla\phi\|_2^2.
 \label{fd-hf-dielectric}
\end{equation}
This follows from the preceding one-particle proof with $\tau=\theta$;
it holds in the mean-charge class and hence also on exact-charge Slaters.

For clarity, one sufficient hypothesis for \eqref{fd-hf-hypothesis} in
the fully free-Wick-ordered functional is the periodic exchange inequality
\[
 \mathrm{Ex}(Q)\le\frac\pi2\operatorname{Tr}(|p|Q^2)
                      +\frac{c_*}{L}\operatorname{Tr}Q^2.
\]
If this inequality has been proved for the selected periodic Green function,
then its nonnegative direct term and \eqref{fd-q-square} give
\[
 \theta=1-\frac{\alpha\pi}{4}-\frac{\alpha c_*}{2mL_0}>0
 \quad\text{uniformly for }L\ge L_0.
\]
This states the needed mean-field assumption explicitly.  The free operator
bound \eqref{fd-dielectric} requires neither that exchange estimate nor
any interacting-vacuum hypothesis.
